\documentclass{article}

% Load the same ICLR style file as the main paper (update year/filename to match uploaded files)
\usepackage{iclr2027_conference}

\usepackage{times}
\usepackage{amsmath, amssymb, amsthm}
\usepackage{booktabs}
\usepackage{graphicx}
\usepackage{natbib}
\usepackage{hyperref}
\usepackage{longtable}
\usepackage{float}
\usepackage{siunitx}
\usepackage{caption}

% Theorem environments for the supplement (consistent with the main paper)
\newtheorem{proposition}{Proposition}
\newtheorem{remark}{Remark}

% Table caption formatting (optional S-prefix labels such as Table S1)
\captionsetup[table]{font=small, skip=4pt}

\title{Beyond High Prices: Auditing Information-Mediated Enforcement in Recurrent Pricing Agents \\ \large Supplementary Material}

\author{Anonymous}

\begin{document}

\maketitle

\appendix

\section*{How to Read This Supplement}

The main paper gives the short argument; this supplement supplies the definitions, derivations, numerical records, and interpretation needed to inspect each step. Sections~A--H retain the technical record and complete probe tables. Section~I explains the main argument in depth, including worked examples, and Section~J gives the public repository address at the end. ``Replaced'' in the main paper means the action-independent ablation recorded as \texttt{ppo\_ablated}; B and A in tables denote informative baseline and ablation. Mathematical firm indices $i\in\{1,2\}$ correspond to archive labels 0 and 1.

\begin{center}
\begin{tabular}{@{}p{0.25\linewidth}p{0.68\linewidth}@{}}
\toprule
Main-paper topic & Where to find the complete explanation \\
\midrule
Question and economic logic & A.5 and I.1 explain earlier work and the difference between prices and enforcement. \\
Market and public reports & B and I.2--I.3 derive benchmarks and trace a cut through aggregation, clipping, and the queue. \\
GRU and PPO & C and I.4 give the architecture, action units, training/deployment distinction, and configuration. \\
Counterfactual audit & A.4, E, and I.5--I.6 define pairing, gains, timing, and the exact negative control. \\
Experimental evidence & D--F and I.7/I.10 provide every seed/probe, uncertainty rules, and worked interpretations. \\
Trigger/noise benchmark & A.1--A.3 and I.8--I.9 give Bellman algebra, quantifiers, derivatives, and limits. \\
Welfare and scope & G--H and I.11--I.12 separate definitions, available evidence, and missing replication inputs. \\
\bottomrule
\end{tabular}
\end{center}


%% ============================================================
\section{Extended Incentive Analysis}
%% ============================================================

This supplement gives extended proofs, the executed PPO configuration, and all 160 conditional price-cut estimates. It reanalyzes the archived ten-seed sample in each condition; no additional policies were trained. The main manuscript frames the study as a causal audit of information-channel dependence, not as a replication of every Q-learning result in the algorithmic-collusion literature.

\subsection{The Value Difference and the Threshold}

Consider the main text's stationary repeated game with absorbing, credible punishment. Write bounded cooperative, punishment, and specified one-period deviation payoffs as $\pi^C$, $\pi^P$, and $\pi^D$. Let $g = \pi^D - \pi^C$ and $\Delta V = \pi^C - \pi^P$. Assume $g > 0$, $\Delta V > 0$, $0\leq\alpha<\beta\leq1$, and $0\leq\delta<1$. The signal depends on current actions; the deviator subsequently follows the prescribed automaton.

Subtract the punishment Bellman equation from the cooperative equation. Both values are evaluated before the current action. This timing gives
\[
V^C - V^P = \frac{\pi^C - \pi^P}{1 - \delta(1 - \alpha)}.
\]
Since $\delta < 1$, the unique value difference is well defined. The deviation changes today's profit and the probability of entering punishment. Therefore
\[
V^D - V^C = g - \delta(\beta - \alpha)(V^C - V^P).
\]
Multiplying by the positive denominator preserves the inequality. The deviation is unprofitable if and only if
\[
\delta \geq \frac{g}{g(1 - \alpha) + (\beta - \alpha)\Delta V}.
\]
The denominator is positive under $\beta>\alpha$, so division yields $\delta^*$. If instead $\beta\leq\alpha$ with $g>0$, the undivided incentive inequality cannot hold and the deviation is always profitable. The strict feasibility condition $\delta^* < 1$ is equivalent to $(\beta - \alpha)\Delta V > \alpha g$. Equality permits only $\delta = 1$, excluded here. If the threshold exceeds one, no admissible discount factor deters the specified deviation. A numerical threshold above one is therefore an infeasibility indicator, not an economically valid patience parameter.

\subsection{Why the Quantifier Matters}

For each firm and action, compute its own current gain and detection probability. It is not enough to test the highest current-profit action. For example, take $\Delta V = 1$ and $\alpha = 0.05$ and consider two deviations. The larger current-gain action, with $g = 0.9$ and $\beta - \alpha = 0.5$, has threshold $180/271 \approx 0.664207$. A smaller current-gain action, with $g = 0.3$ and $\beta - \alpha = 0.01$, has threshold $60/59 \approx 1.016949$. At discount $0.90$, the largest current-gain action is deterred but the smaller one is not, so testing only the former would miss a profitable deviation. These values are an algebraic illustration, not experimental estimates.

The main text's all-action incentive condition must hold for every firm's deviation in $C$, with credible punishment after every history entering $P$. Neither this example nor the finite PPO probes establishes those requirements for all prices and histories.

\subsection{Gaussian Noise and Its Boundaries}

For a fixed threshold $\tau$ between $\mu_D$ and $\mu_C$, let $\sigma$ be the signal standard deviation. Define $z_C = (\tau - \mu_C)/\sigma < 0$ and $z_D = (\tau - \mu_D)/\sigma > 0$. The normal tail probabilities are $\alpha = \Phi(z_C)$ and $\beta = \Phi(z_D)$. Differentiation gives
\[
\frac{\partial \Delta}{\partial \sigma} = -\frac{1}{\sigma}\bigl[z_D \phi(z_D) - z_C \phi(z_C)\bigr] < 0.
\]
Combining these derivatives with the derivatives of $\delta^*$ gives
\[
\frac{\partial \delta^*}{\partial \sigma} = \frac{g\bigl[(g + \Delta V)\alpha'(\sigma) - \Delta V\,\beta'(\sigma)\bigr]}{\bigl[g(1-\alpha) + (\beta-\alpha)\Delta V\bigr]^2} > 0.
\]
Both channels of deterioration matter: more false punishment reduces the cooperative value, and lower detection reduces the expected loss caused by a deviation. As $\sigma$ tends to zero, the threshold tends to $g/(g + \Delta V)$. As $\sigma$ tends to infinity, it tends to $2$. If an additional independent Gaussian noise scale is $\sigma_a$, then $\sigma_{\mathrm{tot}}^2 = \sigma^2 + \sigma_a^2$ and
\[
\frac{\partial \delta^*}{\partial \sigma_a} = \frac{\partial \delta^*}{\partial \sigma_{\mathrm{tot}}} \cdot \frac{\sigma_a}{\sigma_{\mathrm{tot}}}.
\]
For positive base noise, this derivative is positive at $\sigma_a > 0$ and zero at $\sigma_a = 0$. The threshold is still strictly higher at any positive added scale than at zero. This distinction prevents an incorrect claim of a strictly positive derivative at the boundary. The result holds the threshold, means, payoffs, and trigger strategy fixed. It does not sign the response of retrained PPO policies.

\subsection{Observation Invariance and Finite Audit Horizons}

Under the main text's action-independent channel, identical innovations and initial report queues produce identical reports even after an action intervention. Frozen actors receive only report histories and time. Their next prescribed actions therefore coincide; induction establishes equality at every later date. In a stationary payoff model, all later profit differences vanish. This proves the exact negative control without assuming any optimizer convergence.

The result need not hold if actors received action-dependent sales, own-action memory, or rewards as inputs, or if the intervention changed a persistent economic state. During online learning, rewards could also change weights. The archived audit freezes weights, and the PPO execution inputs exclude those additional paths.

For a shorter audit, let $H$ observed post-intervention periods remain within an episode ending at $T$, with periods indexed $0,\ldots,T-1$ and intervention at $t_0$. If both branches' rewards lie in $[0, B]$, the absolute omitted discounted gain is bounded by
\[
\left| \sum_{t=t_0+H+1}^{T-1} \delta^{t-t_0} (\pi_t^{\mathrm{cut}} - \pi_t^{\mathrm{ctrl}}) \right| \leq \frac{\delta^{H+1} B}{1-\delta}.
\]
Here $B = 1.5$ is valid because price minus cost is at most $1.5$ and demand is at most one. The archived probes use $H = T - t_0 - 1$, so this omitted finite-episode tail is zero. This does not bound deviations outside the probe family or establish immunity in an infinite-horizon game.

\subsection{Comparison with Classical and Algorithmic Imperfect Public Monitoring}

Table~\ref{tab:lit_comparison} summarizes how the present study relates to classical imperfect-monitoring repeated-game theory and prior algorithmic pricing studies. Classical repeated-game models \citep{green1984noncooperative, abreu1990toward, fudenberg1994folk} analyze infinite-horizon games in which firms observe a contemporaneous noisy public signal and follow a pre-specified public strategy (such as a two-state threshold trigger or self-generating continuation-value automaton). In those models, dynamic enforcement is characterized analytically by comparing immediate deviation gains against discounted continuation value losses.

In the algorithmic pricing literature, \citet{calvano2020artificial} demonstrate that independent tabular Q-learning agents can converge to supra-competitive prices under perfect price observation, and \citet{calvano2021algorithmic} extend tabular Q-learning to an imperfect public monitoring environment adapted from \citet{green1984noncooperative}. In \citet{calvano2021algorithmic}, firms choose discrete quantities and observe the previous market price under stochastic demand. Their intervention forces an output expansion before restoring learned play. Conversely, recent deep reinforcement learning studies such as \citet{deng2025exploring} compare tabular Q-learning, DQN, and PPO under perfect monitoring and focus primarily on converged price elevation.

\begin{table}[t]
\centering
\caption{Comparison across imperfect public monitoring and algorithmic pricing studies.}
\label{tab:lit_comparison}
\footnotesize
\begingroup
\setlength{\tabcolsep}{3pt}
\begin{tabular}{@{}p{0.20\linewidth}p{0.17\linewidth}p{0.22\linewidth}p{0.11\linewidth}p{0.20\linewidth}@{}}
\toprule
Study & Strategy / Learner & Action \& Signal Space & Horizon & Primary Diagnostic \\
\midrule
\citet{green1984noncooperative}; \citet{abreu1990toward}; \citet{fudenberg1994folk}
  & Fixed trigger / PPE automaton (analytical)
  & Continuous or finite actions; contemporaneous public signal
  & Infinite ($\delta < 1$)
  & Analytical IC threshold \& self-generating continuation sets \\
\addlinespace[3pt]
\citet{calvano2020artificial}
  & Tabular Q-learning
  & Discrete price grid; perfect price observation
  & Infinite
  & Converged markup $\Delta$ \& deviation impulse path \\
\addlinespace[3pt]
\citet{calvano2021algorithmic}
  & Tabular Q-learning with finite memory
  & Discrete quantities; stochastic public market price
  & Infinite
  & Normalized profit \& forced-deviation path \\
\addlinespace[3pt]
\citet{deng2025exploring}
  & TQL, DQN, PPO
  & Discrete prices (including PPO); price histories
  & Discounted; finite training
  & Cross-algorithm price elevation comparison \\
\addlinespace[3pt]
Present study
  & Recurrent PPO (32-step GRU)
  & Continuous prices $[1,2.5]$; delayed 2-period Gaussian report + ablated control
  & Finite $T=512$ with time input
  & Immediate vs.\ continuation payoff decomposition \& proved negative control \\
\bottomrule
\end{tabular}
\endgroup
\end{table}

Our study addresses a distinct methodological and diagnostic gap. Rather than asking only whether converged prices exceed the static Nash benchmark or whether visual price dips follow an exogenous shock, we audit frozen continuous-action neural policies by decomposing the total payoff change from a forced unilateral price cut into its \emph{immediate impact component} ($\Delta\pi_{t_0}$) and its \emph{discounted continuation component} ($\Delta\pi^{\mathrm{cont}}$), paired with an exact action-independent channel replacement (\texttt{ppo\_ablated}, main-text Proposition~1) that provides a proved zero-response negative control. Our findings in the recurrent PPO setting do not refute the tabular Q-learning price-war results of \citet{calvano2021algorithmic}; rather, they show how counterfactual payoff decomposition and channel replacement separate genuine information-mediated continuation threats from immediate static profit losses.

%% ============================================================
\section{Monitoring Calibration and Signal Timing}
%% ============================================================

\subsection{Reference Prices and Parameterization of the Calibration Model}

In the stationary logit market ($c=1$, $v=2$, outside utility $0$, market size $1$, $\mu=0.25$, $p_{i,t}\in[1,2.5]$), the archived elevations use the $1{,}501$-point grid references $p^{\mathrm{N}}=1.473$ and $p^{\mathrm{J}}=1.925$. Independent continuous first-order conditions yield $p^{\mathrm{N}}_{\mathrm{cont}}=1.4729266600$ and $p^{\mathrm{J}}_{\mathrm{cont}}=1.9249809190$. At symmetric prices, these conditions are $1-(p-c)(1-q_i)/\mu=0$ for Nash and $1-(p-c)(1-q_1-q_2)/\mu=0$ for joint profit. They check the static benchmarks, not equilibrium of the learned policies; continuous roots are not substituted into archived elevation calculations. Normalized price elevation $E = (\bar{p}-p^{\mathrm{N}})/(p^{\mathrm{J}}-p^{\mathrm{N}})$ is retained without clipping to $[0,1]$; it is a linear position index between static reference prices, not a structural probability of collusion.

The stated calibration model has six unconstrained coordinates $\theta$. At the fixed threshold $\tau$, it decodes them as
\[
\mu_C = \tau + \exp(\theta_1), \quad \mu_D = \tau - \exp(\theta_2), \quad \sigma = \exp(\theta_3).
\]
Independent normal priors on these coordinates have means $(0.5, 0.5, -0.5, 0, 0, 0)$ and standard deviations $(1, 1, 1, 2, 2, 2)$. Specifically, Gaussian priors are placed on the log distances of the cooperative and deviating means from the threshold $\tau$ and on the log signal scale. The documented numerical support rejects absolute values above $20$ for the log coordinates and above $100$ for the payoff location; thus the implemented target has numerical support bounds. The priors are defined on the coordinates actually sampled, and no change-of-variables Jacobian is added to that coordinate-space target.

Labeled cooperative and deviating signals are modeled with normal likelihoods having means $\mu_C$ and $\mu_D$ and shared standard deviation $\sigma$. Labeled benchmark payoffs also use normal likelihoods with a supplied known payoff observation scale. These labels and restrictions are model inputs; the posterior does not discover a cartel label from an unlabeled price series.

The stated calibration supplies monitoring coordinates $(\theta_{\mathrm{lo}},\theta_{\mathrm{hi}},\sigma)=(\mu_D,\mu_C,\sigma)$ to the PPO environment. The manifest also records \texttt{added\_sigma}=0.1. Because the simulator implementation is absent, its rule combining this setting with posterior $\sigma$ into effective report scale $\sigma_\varepsilon$ remains unverified; we do not equate the two scales or assume a variance-addition rule. Their joint rows are sampled uniformly with replacement at each episode reset, preserving empirical cross-coordinate dependence, and held fixed throughout the episode. Thus, posterior randomization in the PPO experiments represents monitoring uncertainty only; it does not randomize demand or cost primitives. Moreover, the linear interpolation in the main text's latent monitoring rule mapping $\min(p_{1,t},p_{2,t})$ into $[\theta_{\mathrm{lo}},\theta_{\mathrm{hi}}]$ is imposed model structure rather than an empirically estimated relationship between market prices and public reports. Because the latent monitoring map clamps $m_t$ to $[\theta_{\mathrm{lo}},\theta_{\mathrm{hi}}]$, the latent signal mean saturates whenever $\min(p_{1,t},p_{2,t}) \leq p^{\mathrm{N}}$ or $\min(p_{1,t},p_{2,t}) \geq p^{\mathrm{J}}$. Consequently, when baseline prices are near $p^{\mathrm{N}}$, different downward price cuts can generate identical clamped latent means $m_t = \theta_{\mathrm{lo}}$, limiting the channel's sensitivity to cut severity.

\subsection{Timing, Normalization, and Channel Replacement in the Executed Environment}

\begin{table}[t]
\centering
\caption{Signal generation and delivery sequence within one market period.}
\label{tab:timing}
\small
\begin{tabular}{@{}cl@{}}
\toprule
Stage & Operation \\
\midrule
1 & Both firms act using the existing 32-report window and time. \\
2 & The C++ core maps actions to prices and computes demand and profits. \\
3 & The minimum price determines the clipped latent monitoring mean. \\
4 & The last two available latent means are averaged; one normal innovation is added. \\
5 & The generated report enters a one-element queue; the oldest report is delivered. \\
6 & The delivered report is normalized and appended to the next observation. \\
\bottomrule
\end{tabular}
\end{table}

A price change at time $t$ first changes the generated report at $t$. With delay one, that report is delivered after the action at $t+1$, so the earliest actor response is at $t+2$. This accounts for the initial zero entries of the informative impulse response. Overlapping two-period latent-mean windows can spread the signal effect across reports. Normalization of the delivered report is public and fixed across all runs: the delivered report is centered at $1.6$, divided by $0.4$, and clipped to $[-5,5]$. Initial 32-report histories are filled with zeros, and pending queue entries at reset are neutral. No hidden posterior parameter $\theta$ is ever supplied directly to the actor.

The ablation replaces the generated report by $r_{t+1} \sim \mathcal{N}(1.6, 0.4^2)$ independently of prices. Delay, observation window length, economic primitives, and training budgets are unchanged. This replacement enforces exact action independence rather than a numerical approximation to infinite noise variance. However, it also changes the marginal mean and scale of the delivered reports relative to the informative channel. Consequently, retraining under the ablated channel identifies the total effect of this channel replacement, not a pure information-content effect holding all marginal input moments fixed. At the same time, the ablation retains shared randomness: both firms observe the same public draw $r_{t+1}$, and identical innovation consumption across paired branches maintains common-random-number alignment.

%% ============================================================
\section{PPO Implementation and Executed Configuration}
%% ============================================================

The two firms use distinct actor--critic parameters and independent Adam optimizers \citep{schulman2017proximal}. Both firms learn simultaneously in the same market, so single-agent PPO convergence arguments do not imply convergence to a game-theoretic equilibrium. Actors receive neither rivals' prices nor their own realized profits as execution inputs, although each firm's realized rewards train its policy during optimization.

The GRU architecture choice follows directly from the observation structure rather than a claim about commercial pricing software: because the public report averages two consecutive latent price states, adds Gaussian innovation noise, and arrives after a one-step queue delay, a single scalar observation cannot determine whether an adverse report reflects an isolated noise draw or a rival's unilateral price cut \citep{green1984noncooperative, calvano2021algorithmic}. A gated recurrent unit \citep{cho2014gru} provides a compact, differentiable finite-memory sequence encoder over the 32-report window while PPO handles continuous price actions in $[1,2.5]$. Each forward pass recomputes the 128-unit GRU from a fresh zero hidden state across the 32-report window, concatenates normalized episode time $t/512$, and applies a linear projection, $\tanh$, and layer normalization. Thus the implementation is a bounded 32-step sliding-window encoder rather than an unbounded episode-long recurrent state. Policy and value heads share this encoder within each firm, with no weight sharing across firms and no centralized critic. We use GRU+PPO strictly as a controlled testbed for auditing history-dependent continuous pricing policies, and no comparison with an MLP or CNN on the same window is claimed here.

During training, the actor samples a latent Gaussian action and applies $\tanh$ before linearly mapping $[-1,1]$ to prices in $[1,2.5]$ (with action-to-price slope $0.75$). Its log-density includes the analytic $\tanh$ Jacobian correction. This Gaussian action distribution governs training exploration and is distinct from the Gaussian observation noise $\varepsilon_t$, the Gaussian calibration likelihood, and the Gaussian calibration priors (none of which involves Gaussian-process regression or Bayesian optimization). All reported diagnostic curves, final evaluations, and counterfactual audits deploy the frozen deterministic policy $\mu(s)$ obtained by applying $\tanh$ to the latent actor mean. Because $\tanh$ is nonlinear, $\mu(s)$ does not generally equal the expected action under the stochastic training policy; all audit claims apply specifically to this deterministic deployment rule.

\begin{table}[t]
\centering
\caption{Configuration recorded in the completed job files.}
\label{tab:config}
\small
\begin{tabular}{@{}ll@{}}
\toprule
Parameter & Executed Value \\
\midrule
Training seeds in each condition & 20261001 through 20261010 \\
Training interactions per run & \num{1000000} joint market transitions \\
Parallel environments and episode length & 2 environments; 512 periods \\
Public history and encoder width & 32 reports; 128 GRU units \\
Full rollout size & \num{2048} periods per environment; \num{4096} transitions \\
Optimization epochs and minibatch & 10; 256 \\
Adam learning rate and epsilon & \num{0.0003}; \num{0.00001} \\
Discount and GAE parameter & 0.99; 0.95 \\
PPO clipping and entropy coefficient & 0.2; 0.01 \\
Value coefficient and gradient norm cap & 0.5; 0.5 \\
Initial actor log std.\ dev. & $-0.5$; clamped to $[-5, 2]$ during use \\
Informative channel & Delay 1; aggregation 2; added noise scale 0.1 \\
Action-independent channel & Same delay; report $\sim \mathcal{N}(1.6, 0.4^2)$ \\
Diagnostics and final evaluation & 100 checkpoints $\times$ 16 episodes; final 64 episodes \\
Probe family and paired repetitions & 2 firms $\times$ 2 times $\times$ 2 cuts; 128 pairs per probe \\
\bottomrule
\end{tabular}
\end{table}

Advantages are standardized within the complete rollout separately for each firm, using the population standard deviation plus $\epsilon$. For likelihood ratio $r_t$ and standardized advantage $\hat{A}_t$, the minimized loss is
\[
L = -\min(r_t \hat{A}_t, \operatorname{clip}(r_t, 1-\varepsilon, 1+\varepsilon)\hat{A}_t) + c_v \cdot \tfrac{1}{2} \mathrm{MSE} - c_e H.
\]
The effective MSE coefficient is $0.25$ because the code multiplies MSE by $0.5$ and then applies value coefficient $0.5$. Entropy is estimated for the bounded $\tanh$ policy, not its base normal. Stored pre-$\tanh$ actions avoid inverse-$\tanh$ instability.

Episode termination has zero bootstrap; nonterminal rollout boundaries use the next-state value. A shortened final rollout reaches exactly one million interactions. Environments reset synchronously. Deployment uses $\tanh$ of the latent mean; stochastic-policy performance is not reported.

%% ============================================================
\section{Complete Final Outcomes by Training Seed}
%% ============================================================

\begin{table}[H]
\centering
\caption{Final evaluation of every completed PPO policy pair. Each row averages 64 evaluation episodes. B denotes informative baseline and A denotes action-independent retraining. Seed suffix 01 denotes 20261001, continuing through 10. Rank columns describe projected, time-augmented encoder features.}
\label{tab:seed_results}
\small
\sisetup{table-format=1.6, round-mode=places, round-precision=6}
\begin{tabular}{@{}ccS S[table-format=-1.6] SSS@{}}
\toprule
Seed & Arm & {Price} & {Elevation} & {Welfare index} & {Rank 0} & {Rank 1} \\
\midrule
01 & B & 1.497102 & 0.053323 & 1.157867 & 1.276308 & 1.237626 \\
01 & A & 1.488946 & 0.035278 & 1.159020 & 1.418177 & 1.386605 \\
02 & B & 1.481101 & 0.017922 & 1.159558 & 1.218294 & 1.241206 \\
02 & A & 1.476651 & 0.008077 & 1.160178 & 1.547627 & 1.522354 \\
03 & B & 1.468815 & -0.009260 & 1.161126 & 1.229325 & 1.311140 \\
03 & A & 1.458835 & -0.031339 & 1.162020 & 1.328914 & 1.342101 \\
04 & B & 1.479467 & 0.014307 & 1.160148 & 1.254956 & 1.291669 \\
04 & A & 1.477020 & 0.008893 & 1.160382 & 1.232180 & 1.275160 \\
05 & B & 1.509558 & 0.080881 & 1.155365 & 1.281766 & 1.237738 \\
05 & A & 1.466939 & -0.013408 & 1.161483 & 1.331038 & 1.422554 \\
06 & B & 1.461530 & -0.025377 & 1.160801 & 1.259499 & 1.229360 \\
06 & A & 1.487097 & 0.031188 & 1.159354 & 1.293448 & 1.355741 \\
07 & B & 1.483384 & 0.022973 & 1.159627 & 1.265380 & 1.228610 \\
07 & A & 1.493508 & 0.045372 & 1.158365 & 1.280998 & 1.314625 \\
08 & B & 1.463500 & -0.021018 & 1.161653 & 1.342083 & 1.293155 \\
08 & A & 1.493939 & 0.046324 & 1.157274 & 1.310864 & 1.304309 \\
09 & B & 1.474509 & 0.003338 & 1.159803 & 1.280268 & 1.268246 \\
09 & A & 1.491027 & 0.039883 & 1.158652 & 1.291892 & 1.254833 \\
10 & B & 1.479070 & 0.013430 & 1.159930 & 1.335473 & 1.356627 \\
10 & A & 1.480958 & 0.017607 & 1.159871 & 1.351331 & 1.357494 \\
\bottomrule
\end{tabular}
\end{table}

Between-seed price standard deviations are $0.014713$ for B and $0.011766$ for A. Elevation standard deviations are $0.032551$ and $0.026032$. Each run records 100 diagnostic evaluations with 16 episodes each; the requested $10{,}000$-interaction checkpoint spacing is rounded to rollout boundaries, starting at $12{,}288$ interactions. Final outcomes use 64 separate evaluation episodes, and both conditions share seed identifiers and evaluation-stream construction. A mean near the Nash reference does not rule out asymmetry, variation across time, or isolated profitable price cuts. Negative elevation is retained as a valid measured outcome rather than truncated to zero.

Effective feature rank is defined as the exponential Shannon entropy of the normalized eigenvalues of the centered feature covariance matrix. These features are the post-projection, time-augmented representations after $\tanh$ and layer normalization rather than raw GRU hidden states. The reported effective-rank columns are descriptive diagnostics: they do not measure monitoring quality, whether the two firms learned identical representations, or whether a retaliatory punishment rule was encoded. Although the action-independent ablation raises effective rank (from $1.274335$ to $1.338647$ for firm 0 and from $1.269538$ to $1.353578$ for firm 1), it simultaneously changes the marginal input distribution of the public reports. Neither harmful representation collapse under informative monitoring nor a causal effect of feature rank on pricing follows from this comparison.

All 20 rows are included in the reported comparison. A final 64-episode estimate differs slightly from the last 16-episode learning-curve estimate because their evaluation streams differ. The two quantities should not be silently substituted for one another.

%% ============================================================
\section{Statistical Estimands and Uncertainty}
%% ============================================================

Let $Y_s^B$ and $Y_s^A$ denote final outcomes for paired training seed $s$. We estimate the mean channel-replacement effect as
\[
\hat{\Delta} = \frac{1}{10}\sum_{s=1}^{10}(Y_s^A - Y_s^B).
\]
Each bootstrap draw resamples ten paired seed indices with replacement. We report the 2.5th and 97.5th percentiles from \num{2000} draws using NumPy \texttt{default\_rng}, seed 20260920. The same $2000\times10$ index matrix is reused for each outcome, curve, and seed-level audit summary. Whole-seed resampling also gives pointwise learning-curve bands. Intervals condition on the posterior, simulator, hyperparameters, and evaluation protocol; they are not simultaneous bands.

For the counterfactual price-cut audits, we intervene on either firm ($0$ or $1$) at period $t_0\in\{64, 256\}$ using normalized action cuts of $0.1$ or $0.3$, giving eight probe types per seed. Because the action-to-price mapping has slope $0.75$, these normalized cuts correspond to unclipped price reductions of $0.075$ and $0.225$. Each probe uses 128 paired rollouts sharing the episode's posterior draw $\theta$ and Gaussian innovations $\varepsilon_t$; after the single forced cut at $t_0$, both agents return to their frozen deterministic policies through period $511$. We record both the signed sum of the rival impulse response $\Delta p_t^{\mathrm{rival}}$ after impact and the sum of the negative parts of its mean path (a descriptive downward-response area, not a structural measure of punishment). Similarly, the discounted continuation gain $\Delta\pi^{\mathrm{cont}}$ includes both firms' subsequent price adjustments and does not isolate the rival's causal contribution alone. In seed-level summaries, we average probe types within each seed before resampling seeds: cut-specific means average four probe types per seed, while overall rival-response area averages all eight. Seeds receive equal weight; the 128 rollout repetitions and four within-cut probe types are not independent training replications.

For clarity, with $T=512$ and $d_t=\pi_t^{\mathrm{cut}}-\pi_t^{\mathrm{ctrl}}$, the audited full gain is
\[
G=d_{t_0}+\Delta\pi^{\mathrm{cont}},\qquad
\Delta\pi^{\mathrm{cont}}=\sum_{t=t_0+1}^{T-1}\delta^{t-t_0}d_t.
\]
The signed response area is $A=\sum_{t=t_0+1}^{T-1}\Delta p_t^{\mathrm{rival}}$, without discounting. Averaging all eight probes within each informative seed and then resampling seeds gives $\bar A=-1.3628785\times10^{-5}$ and a 95\% interval $[-2.5298314\times10^{-5},-3.4169651\times10^{-6}]$. The small downward response is statistically resolved; it is not itself evidence of continuation deterrence.

\begin{table}[t]
\centering
\caption{Seed-bootstrap intervals for cut-specific full gain and its continuation component. Intervals are pointwise and unadjusted.}
\label{tab:bootstrap_ci}
\small
\sisetup{table-format=-1.6}
\begin{tabular}{@{}cc l SSS@{}}
\toprule
Arm & Cut & Quantity & {Mean} & {Lower 95\%} & {Upper 95\%} \\
\midrule
B & 0.1 & Full gain      & -0.005192 & -0.005855 & -0.004542 \\
B & 0.1 & Continuation   & -0.000005 & -0.000010 & -0.000001 \\
B & 0.3 & Full gain      & -0.051322 & -0.053372 & -0.049286 \\
B & 0.3 & Continuation   & -0.000005 & -0.000010 & -0.000001 \\
A & 0.1 & Full gain      & -0.004800 & -0.005393 & -0.004343 \\
A & 0.1 & Continuation   &  0.000000 &  0.000000 &  0.000000 \\
A & 0.3 & Full gain      & -0.050273 & -0.052023 & -0.048937 \\
A & 0.3 & Continuation   &  0.000000 &  0.000000 &  0.000000 \\
\bottomrule
\end{tabular}
\end{table}

\subsection{Conditional Intervals in the Complete Probe Tables}

Each individual probe file also records a \num{1000}-resample bootstrap interval for the mean across its 128 paired rollouts. The bootstrap generator is initialized with the audit seed plus \num{54321}. Tables~\ref{tab:probe_B01}--\ref{tab:probe_A03} report those stored mean intervals. Their interpretation is conditional on one trained policy pair. The separately stored rollout quantiles named \texttt{deviation\_gain\_q025} and \texttt{deviation\_gain\_q975} are not confidence limits for a mean and are not used in those tables.

In Tables~\ref{tab:probe_B01}--\ref{tab:probe_A03}, seed suffixes 01--10 identify 20261001--20261010; $t$ is intervention time. Immediate denotes impact gain; $G$ includes the remaining episode. Area sums the rival's mean price response after impact, in price-period units. Zeros are rounded except exact action-independent responses. Lower and upper limits are unadjusted conditional mean intervals.

Within-seed probes share policies and random-stream construction. Fixed diagnostic streams across checkpoints likewise do not create independent observations. The 100 checkpoints cannot be treated as 100 independent training samples.

\subsection{What an Unresolved Difference Means}

The price-elevation interval includes zero and effects of either sign. Without a tested equivalence margin, this does not establish equivalence. The four positive conditional small-cut means are descriptive, not multiplicity-adjusted discoveries. Reporting both averages and individual probes preserves this heterogeneity.

%% ============================================================
\section{Complete Probe Tables}
%% ============================================================

% Helper command for probe tables to reduce repetition
\newcommand{\probetable}[4]{%
\begin{longtable}{@{}cccSSSSr@{}}
\caption{#1 \label{#2}} \\
\toprule
Seed & Firm & $t$ & {Immediate} & {$G$} & {Lower 95\%} & {Upper 95\%} & {Area} \\
\midrule
\endfirsthead
\multicolumn{8}{c}{\tablename\ \thetable{} -- continued} \\
\toprule
Seed & Firm & $t$ & {Immediate} & {$G$} & {Lower 95\%} & {Upper 95\%} & {Area} \\
\midrule
\endhead
\bottomrule
\endfoot
#3
\end{longtable}
\vspace{-1em}
\footnotesize{#4}
}

\sisetup{table-format=-1.6, table-number-alignment=center}

\probetable
{Normalized cut 0.1 under informative baseline. Unclipped price cut 0.075. Each row uses 128 paired rollouts.}
{tab:probe_B01}
{
01 & 0 &  64 &  0.000613 &  0.000573 &  0.000557 &  0.000591 & -9.91e-05 \\
01 & 0 & 256 & -0.001242 & -0.001259 & -0.001284 & -0.001234 & -4.23e-05 \\
01 & 1 &  64 & -0.007589 & -0.007606 & -0.007633 & -0.007581 & -3.63e-05 \\
01 & 1 & 256 & -0.007374 & -0.007381 & -0.007422 & -0.007339 & -1.54e-05 \\
02 & 0 &  64 & -0.008006 & -0.008006 & -0.008023 & -0.007989 & 1.28e-07 \\
02 & 0 & 256 & -0.009356 & -0.009356 & -0.009381 & -0.009328 & 0.00e+00 \\
02 & 1 &  64 & -0.000680 & -0.000680 & -0.000701 & -0.000658 & -1.36e-07 \\
02 & 1 & 256 & -0.001313 & -0.001313 & -0.001346 & -0.001279 & 0.00e+00 \\
03 & 0 &  64 & -0.008061 & -0.008061 & -0.008097 & -0.008025 & 0.00e+00 \\
03 & 0 & 256 & -0.008658 & -0.008658 & -0.008711 & -0.008608 & 0.00e+00 \\
03 & 1 &  64 & -0.001741 & -0.001741 & -0.001772 & -0.001713 & 0.00e+00 \\
03 & 1 & 256 & -0.003536 & -0.003536 & -0.003580 & -0.003492 & 0.00e+00 \\
04 & 0 &  64 & -0.005577 & -0.005598 & -0.005633 & -0.005559 & -5.29e-05 \\
04 & 0 & 256 & -0.005955 & -0.005958 & -0.006023 & -0.005893 & -6.00e-06 \\
04 & 1 &  64 & -0.004645 & -0.004677 & -0.004688 & -0.004665 & -8.08e-05 \\
04 & 1 & 256 & -0.005590 & -0.005594 & -0.005611 & -0.005577 & -9.70e-06 \\
05 & 0 &  64 &  0.003834 &  0.003846 &  0.003797 &  0.003899 & 9.39e-06 \\
05 & 0 & 256 &  0.002990 &  0.002990 &  0.002913 &  0.003067 & 1.04e-07 \\
05 & 1 &  64 & -0.009665 & -0.009695 & -0.009729 & -0.009663 & -7.65e-05 \\
05 & 1 & 256 & -0.010022 & -0.010022 & -0.010074 & -0.009969 & -4.99e-07 \\
06 & 0 &  64 & -0.007118 & -0.007118 & -0.007153 & -0.007084 & 0.00e+00 \\
06 & 0 & 256 & -0.005592 & -0.005592 & -0.005639 & -0.005542 & 0.00e+00 \\
06 & 1 &  64 & -0.007011 & -0.007011 & -0.007035 & -0.006985 & 0.00e+00 \\
06 & 1 & 256 & -0.009308 & -0.009308 & -0.009348 & -0.009268 & 0.00e+00 \\
07 & 0 &  64 & -0.005305 & -0.005306 & -0.005323 & -0.005289 & -2.89e-06 \\
07 & 0 & 256 & -0.006025 & -0.006025 & -0.006062 & -0.005991 & 0.00e+00 \\
07 & 1 &  64 & -0.005606 & -0.005605 & -0.005623 & -0.005586 & 1.56e-06 \\
07 & 1 & 256 & -0.005480 & -0.005480 & -0.005509 & -0.005445 & 0.00e+00 \\
08 & 0 &  64 & -0.006527 & -0.006527 & -0.006552 & -0.006503 & -1.20e-08 \\
08 & 0 & 256 & -0.006315 & -0.006315 & -0.006356 & -0.006275 & 0.00e+00 \\
08 & 1 &  64 & -0.004186 & -0.004186 & -0.004196 & -0.004176 & -3.58e-08 \\
08 & 1 & 256 & -0.006705 & -0.006705 & -0.006722 & -0.006688 & 0.00e+00 \\
09 & 0 &  64 & -0.007498 & -0.007528 & -0.007579 & -0.007475 & -7.13e-05 \\
09 & 0 & 256 & -0.010152 & -0.010152 & -0.010233 & -0.010072 & 0.00e+00 \\
09 & 1 &  64 & -0.000496 & -0.000519 & -0.000553 & -0.000488 & -6.26e-05 \\
09 & 1 & 256 & -0.001464 & -0.001464 & -0.001514 & -0.001412 & 0.00e+00 \\
10 & 0 &  64 & -0.008380 & -0.008380 & -0.008424 & -0.008331 & 0.00e+00 \\
10 & 0 & 256 & -0.008557 & -0.008557 & -0.008640 & -0.008481 & 0.00e+00 \\
10 & 1 &  64 & -0.001241 & -0.001241 & -0.001316 & -0.001172 & 0.00e+00 \\
10 & 1 & 256 & -0.002912 & -0.002912 & -0.003033 & -0.002792 & 0.00e+00 \\
}
{Column definitions and conditional-interval limitations appear in Section~E.1.}

\probetable
{Normalized cut 0.3 under informative baseline. Unclipped price cut 0.225. Each row uses 128 paired rollouts.}
{tab:probe_B03}
{
01 & 0 &  64 & -0.032969 & -0.033009 & -0.033065 & -0.032950 & -9.91e-05 \\
01 & 0 & 256 & -0.038903 & -0.038920 & -0.039002 & -0.038841 & -4.23e-05 \\
01 & 1 &  64 & -0.059502 & -0.059519 & -0.059600 & -0.059442 & -3.63e-05 \\
01 & 1 & 256 & -0.058640 & -0.058647 & -0.058773 & -0.058517 & -1.54e-05 \\
02 & 0 &  64 & -0.060599 & -0.060599 & -0.060651 & -0.060547 & 1.28e-07 \\
02 & 0 & 256 & -0.064548 & -0.064548 & -0.064624 & -0.064463 & 0.00e+00 \\
02 & 1 &  64 & -0.037016 & -0.037016 & -0.037085 & -0.036946 & -1.36e-07 \\
02 & 1 & 256 & -0.038793 & -0.038793 & -0.038902 & -0.038682 & 0.00e+00 \\
03 & 0 &  64 & -0.060643 & -0.060643 & -0.060751 & -0.060532 & 0.00e+00 \\
03 & 0 & 256 & -0.062227 & -0.062227 & -0.062388 & -0.062077 & 0.00e+00 \\
03 & 1 &  64 & -0.040365 & -0.040365 & -0.040462 & -0.040275 & 0.00e+00 \\
03 & 1 & 256 & -0.045899 & -0.045899 & -0.046042 & -0.045763 & 0.00e+00 \\
04 & 0 &  64 & -0.052774 & -0.052795 & -0.052906 & -0.052672 & -5.29e-05 \\
04 & 0 & 256 & -0.053804 & -0.053807 & -0.054007 & -0.053610 & -6.00e-06 \\
04 & 1 &  64 & -0.049789 & -0.049821 & -0.049857 & -0.049781 & -8.08e-05 \\
04 & 1 & 256 & -0.052638 & -0.052642 & -0.052691 & -0.052595 & -9.70e-06 \\
05 & 0 &  64 & -0.022114 & -0.022102 & -0.022272 & -0.021922 & 9.42e-06 \\
05 & 0 & 256 & -0.024806 & -0.024806 & -0.025070 & -0.024542 & 1.04e-07 \\
05 & 1 &  64 & -0.065988 & -0.066019 & -0.066124 & -0.065919 & -7.65e-05 \\
05 & 1 & 256 & -0.066949 & -0.066949 & -0.067109 & -0.066787 & -4.99e-07 \\
06 & 0 &  64 & -0.057165 & -0.057165 & -0.057276 & -0.057060 & 0.00e+00 \\
06 & 0 & 256 & -0.052186 & -0.052186 & -0.052337 & -0.052027 & 0.00e+00 \\
06 & 1 &  64 & -0.056827 & -0.056827 & -0.056905 & -0.056744 & 0.00e+00 \\
06 & 1 & 256 & -0.063946 & -0.063946 & -0.064068 & -0.063820 & 0.00e+00 \\
07 & 0 &  64 & -0.051812 & -0.051813 & -0.051868 & -0.051758 & -2.89e-06 \\
07 & 0 & 256 & -0.054033 & -0.054033 & -0.054148 & -0.053925 & 0.00e+00 \\
07 & 1 &  64 & -0.052774 & -0.052774 & -0.052830 & -0.052714 & 1.56e-06 \\
07 & 1 & 256 & -0.052291 & -0.052291 & -0.052384 & -0.052181 & 0.00e+00 \\
08 & 0 &  64 & -0.055728 & -0.055728 & -0.055802 & -0.055653 & -1.20e-08 \\
08 & 0 & 256 & -0.054756 & -0.054756 & -0.054884 & -0.054631 & 0.00e+00 \\
08 & 1 &  64 & -0.048236 & -0.048236 & -0.048270 & -0.048203 & -3.58e-08 \\
08 & 1 & 256 & -0.055997 & -0.055997 & -0.056054 & -0.055942 & 0.00e+00 \\
09 & 0 &  64 & -0.059099 & -0.059129 & -0.059290 & -0.058963 & -7.13e-05 \\
09 & 0 & 256 & -0.066888 & -0.066888 & -0.067128 & -0.066651 & 0.00e+00 \\
09 & 1 &  64 & -0.036522 & -0.036544 & -0.036655 & -0.036442 & -6.26e-05 \\
09 & 1 & 256 & -0.039136 & -0.039136 & -0.039300 & -0.038968 & 0.00e+00 \\
10 & 0 &  64 & -0.061650 & -0.061650 & -0.061788 & -0.061499 & 0.00e+00 \\
10 & 0 & 256 & -0.061991 & -0.061991 & -0.062252 & -0.061754 & 0.00e+00 \\
10 & 1 &  64 & -0.038727 & -0.038727 & -0.038974 & -0.038503 & 0.00e+00 \\
10 & 1 & 256 & -0.043956 & -0.043956 & -0.044347 & -0.043568 & 0.00e+00 \\
}
{Column definitions and conditional-interval limitations appear in Section~E.1.}

\probetable
{Normalized cut 0.1 under action-independent ablation. Unclipped price cut 0.075. Each row uses 128 paired rollouts.}
{tab:probe_A01}
{
01 & 0 &  64 & -0.006554 & -0.006554 & -0.006598 & -0.006508 & 0.00e+00 \\
01 & 0 & 256 & -0.006940 & -0.006940 & -0.007006 & -0.006874 & 0.00e+00 \\
01 & 1 &  64 & -0.001709 & -0.001709 & -0.001766 & -0.001654 & 0.00e+00 \\
01 & 1 & 256 & -0.002856 & -0.002856 & -0.002935 & -0.002775 & 0.00e+00 \\
02 & 0 &  64 & -0.006968 & -0.006968 & -0.007066 & -0.006872 & 0.00e+00 \\
02 & 0 & 256 & -0.008311 & -0.008311 & -0.008476 & -0.008139 & 0.00e+00 \\
02 & 1 &  64 & -0.001395 & -0.001395 & -0.001506 & -0.001278 & 0.00e+00 \\
02 & 1 & 256 & -0.002733 & -0.002733 & -0.002925 & -0.002550 & 0.00e+00 \\
03 & 0 &  64 & -0.002885 & -0.002885 & -0.002902 & -0.002865 & 0.00e+00 \\
03 & 0 & 256 & -0.005016 & -0.005016 & -0.005044 & -0.004988 & 0.00e+00 \\
03 & 1 &  64 & -0.010193 & -0.010193 & -0.010288 & -0.010103 & 0.00e+00 \\
03 & 1 & 256 & -0.009653 & -0.009653 & -0.009781 & -0.009525 & 0.00e+00 \\
04 & 0 &  64 & -0.003106 & -0.003106 & -0.003117 & -0.003095 & 0.00e+00 \\
04 & 0 & 256 & -0.005325 & -0.005325 & -0.005346 & -0.005303 & 0.00e+00 \\
04 & 1 &  64 & -0.005326 & -0.005326 & -0.005348 & -0.005303 & 0.00e+00 \\
04 & 1 & 256 & -0.005506 & -0.005506 & -0.005550 & -0.005464 & 0.00e+00 \\
05 & 0 &  64 & -0.004346 & -0.004346 & -0.004374 & -0.004319 & 0.00e+00 \\
05 & 0 & 256 & -0.005319 & -0.005319 & -0.005352 & -0.005283 & 0.00e+00 \\
05 & 1 &  64 & -0.005744 & -0.005744 & -0.005758 & -0.005730 & 0.00e+00 \\
05 & 1 & 256 & -0.007416 & -0.007416 & -0.007445 & -0.007383 & 0.00e+00 \\
06 & 0 &  64 & -0.003714 & -0.003714 & -0.003759 & -0.003667 & 0.00e+00 \\
06 & 0 & 256 & -0.004358 & -0.004358 & -0.004433 & -0.004278 & 0.00e+00 \\
06 & 1 &  64 & -0.002690 & -0.002690 & -0.002718 & -0.002663 & 0.00e+00 \\
06 & 1 & 256 & -0.004439 & -0.004439 & -0.004478 & -0.004396 & 0.00e+00 \\
07 & 0 &  64 & -0.003969 & -0.003969 & -0.004037 & -0.003901 & 0.00e+00 \\
07 & 0 & 256 & -0.003165 & -0.003165 & -0.003283 & -0.003038 & 0.00e+00 \\
07 & 1 &  64 & -0.003615 & -0.003615 & -0.003674 & -0.003557 & 0.00e+00 \\
07 & 1 & 256 & -0.005815 & -0.005815 & -0.005924 & -0.005725 & 0.00e+00 \\
08 & 0 &  64 & -0.000468 & -0.000468 & -0.000531 & -0.000412 & 0.00e+00 \\
08 & 0 & 256 &  0.000122 &  0.000122 &  0.000030 &  0.000219 & 0.00e+00 \\
08 & 1 &  64 & -0.007390 & -0.007390 & -0.007417 & -0.007362 & 0.00e+00 \\
08 & 1 & 256 & -0.009166 & -0.009166 & -0.009213 & -0.009120 & 0.00e+00 \\
09 & 0 &  64 & -0.007326 & -0.007326 & -0.007358 & -0.007295 & 0.00e+00 \\
09 & 0 & 256 & -0.007760 & -0.007760 & -0.007813 & -0.007705 & 0.00e+00 \\
09 & 1 &  64 & -0.000717 & -0.000717 & -0.000728 & -0.000705 & 0.00e+00 \\
09 & 1 & 256 & -0.001924 & -0.001924 & -0.001942 & -0.001906 & 0.00e+00 \\
10 & 0 &  64 & -0.006716 & -0.006716 & -0.006791 & -0.006631 & 0.00e+00 \\
10 & 0 & 256 & -0.007488 & -0.007488 & -0.007616 & -0.007368 & 0.00e+00 \\
10 & 1 &  64 & -0.001349 & -0.001349 & -0.001423 & -0.001282 & 0.00e+00 \\
10 & 1 & 256 & -0.002766 & -0.002766 & -0.002873 & -0.002654 & 0.00e+00 \\
}
{Column definitions and conditional-interval limitations appear in Section~E.1.}

\probetable
{Normalized cut 0.3 under action-independent ablation. Unclipped price cut 0.225. Each row uses 128 paired rollouts.}
{tab:probe_A03}
{
01 & 0 &  64 & -0.056118 & -0.056118 & -0.056257 & -0.055972 & 0.00e+00 \\
01 & 0 & 256 & -0.057140 & -0.057140 & -0.057347 & -0.056930 & 0.00e+00 \\
01 & 1 &  64 & -0.040510 & -0.040510 & -0.040694 & -0.040329 & 0.00e+00 \\
01 & 1 & 256 & -0.044039 & -0.044039 & -0.044295 & -0.043776 & 0.00e+00 \\
02 & 0 &  64 & -0.057397 & -0.057397 & -0.057704 & -0.057098 & 0.00e+00 \\
02 & 0 & 256 & -0.061263 & -0.061263 & -0.061778 & -0.060726 & 0.00e+00 \\
02 & 1 &  64 & -0.039447 & -0.039447 & -0.039814 & -0.039062 & 0.00e+00 \\
02 & 1 & 256 & -0.043422 & -0.043422 & -0.044051 & -0.042825 & 0.00e+00 \\
03 & 0 &  64 & -0.043619 & -0.043619 & -0.043688 & -0.043543 & 0.00e+00 \\
03 & 0 & 256 & -0.050355 & -0.050355 & -0.050460 & -0.050253 & 0.00e+00 \\
03 & 1 &  64 & -0.066864 & -0.066864 & -0.067150 & -0.066595 & 0.00e+00 \\
03 & 1 & 256 & -0.065033 & -0.065033 & -0.065421 & -0.064651 & 0.00e+00 \\
04 & 0 &  64 & -0.045067 & -0.045067 & -0.045105 & -0.045031 & 0.00e+00 \\
04 & 0 & 256 & -0.051887 & -0.051887 & -0.051956 & -0.051816 & 0.00e+00 \\
04 & 1 &  64 & -0.052212 & -0.052212 & -0.052282 & -0.052141 & 0.00e+00 \\
04 & 1 & 256 & -0.052465 & -0.052465 & -0.052601 & -0.052335 & 0.00e+00 \\
05 & 0 &  64 & -0.048844 & -0.048844 & -0.048928 & -0.048761 & 0.00e+00 \\
05 & 0 & 256 & -0.051607 & -0.051607 & -0.051705 & -0.051502 & 0.00e+00 \\
05 & 1 &  64 & -0.053324 & -0.053324 & -0.053366 & -0.053283 & 0.00e+00 \\
05 & 1 & 256 & -0.058280 & -0.058280 & -0.058365 & -0.058184 & 0.00e+00 \\
06 & 0 &  64 & -0.047331 & -0.047331 & -0.047471 & -0.047183 & 0.00e+00 \\
06 & 0 & 256 & -0.049064 & -0.049064 & -0.049295 & -0.048818 & 0.00e+00 \\
06 & 1 &  64 & -0.044019 & -0.044019 & -0.044108 & -0.043932 & 0.00e+00 \\
06 & 1 & 256 & -0.049324 & -0.049324 & -0.049447 & -0.049195 & 0.00e+00 \\
07 & 0 &  64 & -0.047984 & -0.047984 & -0.048202 & -0.047770 & 0.00e+00 \\
07 & 0 & 256 & -0.045168 & -0.045168 & -0.045549 & -0.044759 & 0.00e+00 \\
07 & 1 &  64 & -0.046842 & -0.046842 & -0.047032 & -0.046656 & 0.00e+00 \\
07 & 1 & 256 & -0.053689 & -0.053689 & -0.054034 & -0.053398 & 0.00e+00 \\
08 & 0 &  64 & -0.036450 & -0.036450 & -0.036653 & -0.036267 & 0.00e+00 \\
08 & 0 & 256 & -0.034251 & -0.034251 & -0.034556 & -0.033933 & 0.00e+00 \\
08 & 1 &  64 & -0.058781 & -0.058781 & -0.058867 & -0.058688 & 0.00e+00 \\
08 & 1 & 256 & -0.064135 & -0.064135 & -0.064281 & -0.063989 & 0.00e+00 \\
09 & 0 &  64 & -0.058559 & -0.058559 & -0.058655 & -0.058464 & 0.00e+00 \\
09 & 0 & 256 & -0.059720 & -0.059720 & -0.059881 & -0.059553 & 0.00e+00 \\
09 & 1 &  64 & -0.037251 & -0.037251 & -0.037291 & -0.037208 & 0.00e+00 \\
09 & 1 & 256 & -0.040991 & -0.040991 & -0.041055 & -0.040926 & 0.00e+00 \\
10 & 0 &  64 & -0.056647 & -0.056647 & -0.056881 & -0.056381 & 0.00e+00 \\
10 & 0 & 256 & -0.058798 & -0.058798 & -0.059196 & -0.058423 & 0.00e+00 \\
10 & 1 &  64 & -0.039348 & -0.039348 & -0.039594 & -0.039123 & 0.00e+00 \\
10 & 1 & 256 & -0.043667 & -0.043667 & -0.044018 & -0.043295 & 0.00e+00 \\
}
{Column definitions and conditional-interval limitations appear in Section~E.1.}

%% ============================================================
\section{Reproducibility and Computation Record}
%% ============================================================

\subsection{Provenance of the Reported Outputs}

The results archive contains final PPO checkpoints, job configurations, learning curves, probe summaries, and suite metadata. All 20 final checkpoint file hashes match their job records in the present audit. The manifest records Python-source, native-extension, and posterior identities, but their corresponding files are absent. Recorded hashes alone do not verify unavailable source bytes, simulator behavior, calibration, or the validity of a completed run. Implementation details stated in Sections~B--C describe the documented design; they cannot be independently confirmed from these outputs.

The recorded environment uses Python 3.13.15, NumPy 2.1.3, PyTorch 2.11.0+cu128, Gymnasium 1.3.0, and device \texttt{cuda}. These are archived versions, not a new compatibility guarantee.

\begin{table}[H]
\centering
\caption{Recorded identities of unavailable replication inputs; these hashes alone do not verify their contents.}
\label{tab:hashes}
\small
\begin{tabular}{@{}ll@{}}
\toprule
Component & SHA256 \\
\midrule
Posterior draws & \texttt{a2cd4f268b3edf0d...202d1b95} \\
Native extension & \texttt{8c1fe82a4830d09e...ef882780} \\
PPO training source & \texttt{12ffcf180b6fc899...662d3f15} \\
\bottomrule
\end{tabular}
\end{table}

The diagnostic gate records maximum rank-normalized split $\hat{R} = 1.008945$ and minimum bulk/tail ESS $= 606.86$, satisfying the software convergence gates of $1.01$ and $400$ \citep{vehtari2021rank}. Passing these thresholds addresses MCMC sampling quality, not Gaussianity of real-world monitoring errors or external market calibration. This supplement does not reconstruct missing likelihood inputs or invent posterior moments from those two statistics. Independent calibration replication requires the original labeled observations, observation-scale settings, sampling configuration, and retained draws matching the recorded posterior hash.

\subsection{Interaction Accounting}

\begin{table}[H]
\centering
\caption{Recorded interaction counts across all completed jobs.}
\label{tab:interactions}
\small
\begin{tabular}{@{}lrr@{}}
\toprule
Component & Per Job & Across 20 Jobs \\
\midrule
Training              & \num{1000000}  & \num{20000000} \\
Training diagnostics  & \num{819200}   & \num{16384000} \\
Trainer final eval.   & \num{8192}     & \num{163840} \\
Suite final eval.     & \num{32768}    & \num{655360} \\
Paired deviation audit & \num{1048576} & \num{20971520} \\
\midrule
Total recorded        & \num{2908736}  & \num{58174720} \\
\bottomrule
\end{tabular}
\end{table}

Each audit type costs \num{131072} transitions, including the pre-intervention prefix in both branches. Eight types give the reported \num{1048576}. One market transition contains both firms' actions; it is not two separate environment interactions. Recorded completed-job times sum to $21.978$ hours for baseline and $22.079$ hours for ablation ($44.06$ hours total). Hardware-specific runtime comparisons require additional telemetry beyond these job totals.

Three previous-attempt records concern baseline seed 20261010 and ablated seeds 20261008 and 20261009. All are marked running, not failed, in the archived records. Their filenames contain \texttt{previous\_failure}, but the records do not establish a numerical training failure or its cause. All three final jobs completed, and unrecorded interrupted work is excluded from the interaction and runtime totals above. Finally, the broader research plan also scheduled Soft Actor-Critic (SAC) conditions, but no incomplete SAC comparison enters this study; because the PPO-only scope and its interpretation were finalized after inspecting completed runs, the manuscript is not a preregistered confirmatory test of a null effect.

%% ============================================================
\section{Detailed Limitations, Reproduction Steps, and Extensions}
%% ============================================================

\subsection{Expanded Discussion of Methodological and Structural Limitations}

Several structural boundaries govern what the analytical benchmark and the PPO audit can and cannot establish:
\begin{enumerate}
\item \textbf{Conditional connection between the infinite-horizon benchmark and finite-horizon PPO:} The analytical trigger benchmark in Section~A assumes an infinite horizon with contemporaneous binary threshold signals, whereas the PPO environment passes continuous, two-period-averaged, one-step-delayed reports to a 32-step GRU over a known finite horizon of $T=512$ periods with normalized time included in the actor input and zero terminal value bootstrap. In a complete-information finite repeated game with a unique stage-game Nash equilibrium, backward induction selects stage-equilibrium play in every period. We do not invoke backward induction as a formal convergence theorem for simultaneous gradient-based neural learners under imperfect monitoring, but we emphasize that a known terminal date and observable remaining time fundamentally alter the dynamic learning problem. The near-Nash outcomes in our PPO runs neither refute the infinite-horizon trigger theorem nor prove that neural learners cannot collude in continuing tasks.
\item \textbf{Input distribution confound:} Replacing the informative channel with $r_{t+1}\sim\mathcal{N}(1.6,0.4^2)$ eliminates action dependence but also changes the marginal mean, variance, and autocorrelation of the report stream. Separating pure information-content effects from input-scale and architecture effects requires moment-matched action-independent placebo channels, history-length controls, and nonrecurrent (MLP) baselines.
\item \textbf{Signal saturation near Nash:} Because the latent monitoring map clamps $m_t$ at $\theta_{\mathrm{lo}}$ when $\min(p_{1,t},p_{2,t})\leq p^{\mathrm{N}}$, price cuts from baseline prices near $p^{\mathrm{N}}$ can saturate the latent signal mean. Nonsaturating signal maps and smaller, adaptive price cuts are needed before interpreting small impulse responses as insensitivity to economically relevant departures.
\item \textbf{Probe coverage and positive deviations:} Our finite probe family tests two downward cut sizes at two intervention periods. Among the 40 small-cut probes per condition, 3 of 40 under informative monitoring (across two seeds, maximum conditional mean gain $0.003846$) and 1 of 40 under action-independent monitoring (gain $0.000122$) exhibit positive conditional mean gains whose within-policy bootstrap intervals exclude zero individually (unadjusted for multiplicity across probes). Testing upward price changes, best-response searches, multistep deviations, and alternative histories could uncover additional profitable deviations.
\item \textbf{Missing collusive positive control:} Proposition~1 proves that our action-independent channel serves as an exact zero-response negative control ($\Delta p_t^{\mathrm{rival}}=0$ and $\Delta\pi^{\mathrm{cont}}=0$), but a negative control does not establish statistical power (sensitivity). Validating the audit as a detector with known sensitivity requires evaluating it on a known collusive positive control---such as an explicitly programmed noisy trigger policy or a verified collusive learner---that generates the predicted tradeoff between positive current deviation gains ($g>0$) and larger discounted continuation losses.
\item \textbf{Model scope and welfare accounting:} All inference is conditional on one retained posterior, one stationary logit demand specification, and ten independent training seeds. A confidence interval spanning zero ($[-0.0241,\,0.0288]$ for price elevation difference) reflects finite-sample uncertainty and is not proof of zero effect. Section~\ref{app:welfare} defines standard logit total surplus as inclusive-value consumer surplus plus profits, with the outside option as normalization. The simpler expression $(v-c)(q_1+q_2)$ omits preference-surplus terms and is bounded by one, so it cannot represent the archived values near $1.1596$. We retain those values as an \emph{archived welfare index} (the field \texttt{mean\_total\_surplus}), without claiming to have verified its simulator formula. Aggregate mean prices and profits cannot reconstruct nonlinear, period-level welfare; that requires the missing simulator source and price trajectories.
\end{enumerate}

\subsection{Reproduce the Reported PPO Scope}

This repository alone cannot execute the training reproduction: it lacks the \texttt{acrl} Python package, native simulator source/binary, build instructions, and calibration input files. The following is a conditional recipe requiring those original materials; it was not executed in this audit. Supply the original \texttt{posterior\_draws.csv} and matching \texttt{posterior\_summary.csv}, verify the posterior hash, and run \texttt{require\_converged}. Preserve joint posterior rows; independent marginal sampling changes the training distribution.

Copy \texttt{experiments/research.json}, retaining \texttt{ppo\_baseline}, \texttt{ppo\_ablated}, and their paired comparison. Keep seeds, evaluation counts, and audit settings unchanged. Save as \texttt{experiments/ppo\_research.json}. The original suite also scheduled SAC; this PPO-only recipe changes plan identity and requires a new output directory.

Run \texttt{python -m acrl.experiment\_suite} with \texttt{-\/-plan experiments/ppo\_research.json}, \texttt{-\/-posterior results/mcmc/posterior\_draws.csv}, \texttt{-\/-output results/ppo\_reproduction}, and \texttt{-\/-device cuda}. CPU changes run identity and may change floating-point behavior. Exact replication across hardware or package versions requires verification.

The \texttt{-\/-resume} option reuses validated completed jobs with matching identities. It does not continue an interrupted optimizer state. A progress policy checkpoint alone does not restore all optimizer and random-generator states. Preserve the prior archive and identify restarted jobs.

Check job statuses, hashes, and finite metrics. Recompute seed summaries from \texttt{job.json}, curves from \texttt{policy.curve.csv}, and probes from \texttt{audit\_probes.json} using Section~E's conventions. Exclude unfinished SAC runs.

\subsection{Extensions That Would Strengthen the Evidence}

The following are not completed experiments. They define limits of the conclusions. A Calvano-style Q-learning replication provides an external benchmark and can serve as an empirical positive control if its enforcement is independently verified. A designed trigger strategy supplies a control with a known current-gain/continuation-loss tradeoff.

\begin{table}[t]
\centering
\caption{Unresolved extensions and the questions they would address.}
\label{tab:extensions}
\small
\begin{tabular}{@{}ll@{}}
\toprule
Additional Experiment & Question It Would Resolve \\
\midrule
Known collusive trigger as positive control & Does the audit detect a designed gain/loss tradeoff? \\
Calvano-style tabular Q-learning benchmark & Does the implementation reproduce prior results? \\
Matched-moment action-independent reports & Is the effect due to information or input scale? \\
Nonsaturating signal maps and smaller cuts & Does saturation near Nash conceal sensitivity? \\
Recurrent versus memoryless policies & Does history processing change enforcement? \\
Adaptive and multistep deviations & Do policies admit profits outside fixed probes? \\
Continuing-task and horizon sensitivity & How much do terminal date and finite memory matter? \\
Gaussian-noise sweep with fixed trigger & Does the comparative static hold empirically? \\
Alternative learners and market primitives & How far do PPO findings generalize? \\
\bottomrule
\end{tabular}
\end{table}

\section{A Detailed Walkthrough of the Main Paper}\label{app:walkthrough}

This section expands the main paper's shorter explanations. It follows the causal chain that the audit is intended to test: a firm's price cut changes the public report, the changed report affects future actions, and those actions change future profits. The examples below interpret existing archived estimates or illustrate the analytical model. They are not additional training runs or newly observed data.

\subsection{The Central Distinction: A High Price Is an Outcome, Not a Mechanism}

A high price describes what a policy chooses. An enforcement mechanism explains why a firm would reject a tempting alternative. The difference matters because a learning algorithm can select a high price for reasons that have nothing to do with punishment. Limited exploration, an imperfect function approximation, or an unfinished optimization process can all produce a price above a static reference. The present experiments do not identify which of these explanations causes early elevation; they show that early elevation alone is insufficient evidence of dynamic enforcement.

The mechanism we seek has three economic parts. First, cutting the current price must offer a current benefit. Second, the rival must obtain information that changes the subsequent action path. Third, the resulting continuation loss must outweigh that current benefit. A downward rival price response addresses only the second part. A negative total gain from the cut also does not establish all three parts, because the first part may already fail. The audit therefore reports both the immediate and continuation components instead of interpreting a single full-gain number as proof of punishment.

The word ``punishment'' can be used more loosely to describe any adverse reaction. Our narrower use concerns \emph{deterrence}: future consequences make an otherwise profitable action unattractive. A policy may react to a report without producing enough of a loss to deter a deviation. Conversely, a policy can fail the tested cut without admitting that every possible deviation is profitable. These distinctions explain the main paper's restricted conclusion: there is no evidence of continuation deterrence for the specified frozen PPO policies and tested one-step cuts.

\subsection{How the Market Creates a Static Tradeoff}

The logit demand expression in the main paper has an outside-option term equal to one in the denominator. A buyer can therefore choose either firm or neither firm, and $q_1+q_2<1$. Lowering $p_i$ raises the firm's demand but lowers its margin $p_i-c$. Differentiating demand with respect to the firm's own price gives
\[
\frac{\partial q_i}{\partial p_i}=-\frac{q_i(1-q_i)}{\mu}.
\]
The corresponding profit derivative is
\[
\frac{\partial\pi_i}{\partial p_i}
=q_i\left[1-\frac{(p_i-c)(1-q_i)}{\mu}\right].
\]
At an interior best response the bracket is zero. This is an individual-firm calculation: the firm considers its own profit with the rival's current price fixed. Under symmetric prices, numerical solution gives the continuous Nash reference $1.4729266600$. The archived normalization instead uses the nearby grid point $1.473$. Both statements can be true; changing between them silently would change every reported elevation.

For the symmetric joint-profit calculation, write $Q=q_1+q_2$ and let both prices vary together. The derivative of $(p-c)Q(p,p)$ has the factor $1-(p-c)(1-Q)/\mu$. Its root is $1.9249809190$, near the archived grid reference $1.925$. The two benchmarks answer different questions: individual incentives versus coordinated total profit. The elevation index is only the linear location of a mean price between these references. It is not the probability that agents collude, and negative values are valid measurements.

These derivatives also explain an apparently puzzling empirical result. A cut below a price near an individual best response can reduce current profit because the margin loss exceeds the demand gain. That loss does not require any future price response. Accordingly, a cut of a fixed size is an intervention to study, not a guarantee of a positive immediate gain. Finding the best unilateral deviation would require searching the action space rather than assuming that either fixed cut is best.

\subsection{Following a Price Cut Through the Report Queue}

The report channel is deliberately indirect. It does not reveal a pair of prices. The minimum price is mapped into a latent mean, two successive means are averaged, noise is added, and the report waits in a queue. An individual report can consequently be low because of a cut, an earlier price, or a noise realization. More history can help a policy filter these possibilities, but history does not by itself guarantee correct inference or an enforcement strategy.

To follow timing, consider a cut at $t_0$. Actions at $t_0$ use the histories already available. The cut can change $m_{t_0}$ and the report generated after that action. The queue still delivers an older report. At $t_0+1$, actions are again selected before the changed report is delivered. Only at $t_0+2$ can an actor use the report containing the cut. This yields zero rival responses at impact and lag one under the documented timing. The archived mean paths have those two zero entries; source-code confirmation of the queue remains unavailable.

Because the average includes both $m_t$ and $m_{t-1}$, one changed latent mean can enter two generated reports. Any later persistence can then arise from the overlapping report windows, the actor's finite memory, or subsequent actions feeding back into reports. It should not automatically be interpreted as the duration of an intended punishment state. A GRU can implement many continuous response rules, not just a two-state automaton.

Clipping creates a separate observation limit. At $\min(p_1,p_2)\le p^{\mathrm N}$, the mapped mean is already $\theta_{\mathrm{lo}}$. Further cuts leave that mean unchanged. The two tested cut sizes can therefore produce similar report changes even though their immediate profit effects differ greatly. A nonsaturating map or smaller, history-adaptive interventions could help distinguish weak policy sensitivity from weak transmitted signals. Those extensions have not been run here.

The effective report-noise scale must also be distinguished from the posterior scale and the recorded added-noise setting. The metadata show \texttt{added\_sigma}=0.1; they do not reveal the missing simulator's combination rule. An independent-Gaussian variance sum is a valid mathematical construction, but treating it as the implemented rule would be an unsupported inference.

\subsection{What the GRU and PPO Do, and What the Audit Freezes}

The GRU processes the 32-report observation window as a sequence. Its hidden state is restarted at zero for each forward pass, so its memory is a sliding window rather than an unrestricted record of the whole episode. Normalized episode time is then combined with the encoded reports. This matters for interpretation: time is an execution input, and a policy can behave differently near the known terminal date even when the report history looks similar.

PPO changes network parameters during training using sampled actions and rewards. A reward can train the policy without being an observation supplied to the deployed actor. The distinction is essential to the negative-control proof. In the frozen audit, the weights stop changing, and actors receive only the stated report histories and time. They do not receive the intervention's realized profit through an additional observation channel.

If $a\in[-1,1]$ denotes normalized action, the linear price mapping is $p=1.75+0.75a$. A forced action cut $u$ replaces $a$ by $\operatorname{clip}(a-u,-1,1)$. Without clipping, price falls by $0.75u$: $0.075$ for $u=0.1$ and $0.225$ for $u=0.3$. The recorded realized cuts can be checked in the probe files. Describing the nominal action cut as a price cut of the same number would mix units.

The stochastic training rule first samples a latent Gaussian action and then applies $\tanh$. Deployment instead applies $\tanh$ to the latent mean. Nonlinearity means that the latter generally differs from the stochastic policy's expected action. We consequently avoid extending a deterministic-deployment audit to all stochastic deployment rules. Likewise, successful single-agent policy optimization does not prove convergence to a Nash or public perfect equilibrium when two policies learn together.

\subsection{Why Pairing Helps, and Exactly What the Gain Measures}

The control and cut branches share the monitoring draw, pre-intervention history, and innovation sequence. Pairing keeps random market reports from obscuring the profit difference caused by the intervention. It does not make the branches' future actions equal under informative reports: the changed price may change reports and therefore policy inputs. Such divergence is precisely the pathway being measured.

The gain is anchored at $t_0$, so the immediate profit difference has weight one and the next period has weight $\delta$. It includes every remaining period through $511$. A negative continuation term means that the deviator earns less after impact in the cut branch. However, both the deviator and rival can change future prices in response to their shared altered reports. The continuation term therefore measures the total future effect of the intervention on the focal firm's profit, not an isolated causal effect of the rival's actions.

The signed rival-response area is a different estimand. It sums price differences without profit weighting or discounting. Positive and negative price movements may offset. The downward-response statistic instead sums negative parts of the mean response path. Taking the negative part after averaging is not the same operation as averaging the negative part of each individual rollout. Neither area is automatically a measure of an incentive constraint. Reporting their definitions avoids treating similarly named response summaries as interchangeable.

The finite-episode tail is zero for these archived probes because all remaining periods are included. This does not mean there is no economically relevant behavior after a continuing-market horizon, or that multistep deviations have been bounded. The finite tail calculation addresses missing periods inside the specified episode only; Section~A.4 states that scope explicitly.

\subsection{A Step-by-Step View of the Negative Control}

Before the intervention, the branches have the same prescribed actions because their inputs and frozen parameters agree. At impact, the intervention externally replaces one selected action; it does not change the parameters or time input. Under the replacement channel, the next generated report depends on a shared random draw, not either branch's prices. The queues and observation windows therefore remain equal.

Now assume the two branches have equal windows before some subsequent action. Their deterministic actors return equal actions. The price-independent report generator then appends equal reports, preserving equality before the next action. Induction gives equality at every later date. Since the market has no persistent payoff-relevant state in the stated design, equal current prices yield equal current profits. Every post-impact profit difference is zero, making the continuation gain exactly zero in the mathematical model.

Numerically, the 80 stored replaced-channel mean response paths are zero throughout. Full and immediate gain summaries can differ at the order of floating-point roundoff because they are aggregated separately. Such a difference is not an economically meaningful violation of the theorem. More importantly, observing zero in stored summaries does not independently prove that unavailable source code implements every assumption.

The counterfactual could change if actors remembered their own executed actions, observed intervention-dependent rewards or sales, updated their weights online, or interacted through a persistent demand state. These are concrete alternative pathways, not technical details that can be omitted from an invariance claim. The control is exact only because the stated audit excludes them.

\subsection{Why Ten Seeds, Not Thousands of Independent Samples}

For a final-outcome comparison, a paired observation is $(Y_s^B,Y_s^A)$ for training seed $s$. The bootstrap draws ten seed indices with replacement, takes the corresponding paired differences, and averages them. Repeating this 2,000 times approximates the sampling distribution of the ten-seed mean. The interval uses the 2.5th and 97.5th percentiles. The fixed index matrix and generator seed are specified in Section~E so that these summaries can be regenerated.

An audit summary follows the same unit-of-replication rule. Average the four firm--time probe types for a given cut within each seed, then resample the ten seed averages. For the overall response area, average all eight types within each seed before resampling. Counting each of the 128 rollout pairs or each diagnostic checkpoint as an independent trained policy would greatly overstate replication. Repeated observations of a policy do not replace independently trained policies.

A conditional probe interval answers a narrower question: given this one trained policy pair, what is the uncertainty in its average over the evaluation rollouts? The archive stores these intervals but not the individual rollout gains needed to regenerate them. Their equality to the printed probe tables is verified; their underlying bootstrap cannot be rerun independently from summaries alone. Stored rollout quantiles describe the spread of individual outcomes and must not be substituted for confidence limits on a mean.

The paired elevation interval spans zero. It allows either sign of an average channel-replacement effect under this protocol. It does not show that the true effect is zero, and no equivalence margin was prespecified or tested. Conversely, the corrected response-area interval excludes zero, so describing the entire informative response as statistically indistinguishable from zero would be wrong. The practical question of enforcement still depends on magnitudes and profits, not only whether one interval crosses zero.

\subsection{Deriving the Trigger Threshold Without Skipping Steps}

Use $h=\Delta V=\pi^C-\pi^P$ for the \emph{period-payoff} gap and $L=V^C-V^P$ for the \emph{continuation-value} gap. These are different objects. In the absorbing punishment state,
\[
V^P=\pi^P+\delta V^P=\frac{\pi^P}{1-\delta}.
\]
Before acting in the cooperative state,
\[
V^C=\pi^C+\delta[(1-\alpha)V^C+\alpha V^P].
\]
Subtracting the first Bellman equation from the second gives
\[
L=h+\delta(1-\alpha)L,
\qquad L=\frac{h}{1-\delta(1-\alpha)}.
\]
The denominator is positive for $\delta<1$. A larger false-alarm probability reduces $L$: cooperation is less valuable if an undeserved punishment can interrupt it sooner.

The specified deviator gains $g$ now, but raises the probability of the lower continuation state from $\alpha$ to $\beta$. Its expected loss of continuation value is therefore $\delta(\beta-\alpha)L$, giving
\[
V^D-V^C=g-\delta(\beta-\alpha)L.
\]
The deviation is unprofitable when the second term is at least $g$. Multiplying by the positive denominator yields
\[
g[1-\delta(1-\alpha)]\le\delta(\beta-\alpha)h.
\]
Rearrangement gives $g\le\delta D$ with $D=g(1-\alpha)+(\beta-\alpha)h$. Only after imposing $\beta>\alpha$, $g>0$, and $h>0$ is it justified to divide by $D>0$ and obtain the displayed threshold. If $\beta\le\alpha$, the expected future loss cannot offset a strictly positive current gain. A threshold below one exists exactly when $D>g$, or $(\beta-\alpha)h>\alpha g$.

The complete incentive condition repeats this comparison for each firm and action. Section~A.2 illustrates why the currently most profitable deviation need not be the hardest one to deter. Punishment credibility is another requirement: prescribed play in $P$ must itself be optimal. Here it is assumed to be credible Nash play. None of these fixed-strategy conditions proves that PPO finds the strategy.

\subsection{The Gaussian Comparison and Its Restrictions}

At a fixed threshold between the two means, set $z_C=(\tau-\mu_C)/\sigma<0$ and $z_D=(\tau-\mu_D)/\sigma>0$. The Gaussian CDF gives $\alpha=\Phi(z_C)$ and $\beta=\Phi(z_D)$. Since the derivative of $z$ with respect to $\sigma$ is $-z/\sigma$,
\[
\alpha'=-z_C\phi(z_C)/\sigma>0,
\qquad \beta'=-z_D\phi(z_D)/\sigma<0.
\]
For $\delta^*=g/D$, differentiating with respect to the probabilities at fixed $g,h$ gives
\[
\frac{\partial\delta^*}{\partial\alpha}=\frac{g(g+h)}{D^2}>0,
\qquad
\frac{\partial\delta^*}{\partial\beta}=-\frac{gh}{D^2}<0.
\]
The chain rule therefore yields a strictly positive derivative with respect to $\sigma$. This is a two-part deterioration: more mistaken punishments and fewer detected deviations. Holding the trigger and payoffs fixed is essential. If agents retrain, means move, or the threshold is changed with noise, these derivatives do not determine the net result.

At very low noise, the threshold separates the two distributions almost perfectly, so $\alpha\to0$ and $\beta\to1$. At very high noise both approach $1/2$, eliminating useful detection differences. The algebraic threshold then approaches two. This does not posit a discount factor of two; it says no admissible discount factor below one can sustain that specified trigger incentive. The same logic explains why action-independent reports provide no action-based continuation deterrence in the restricted frozen-input model.

If an independent extra Gaussian disturbance were added in quadrature, the total scale would satisfy $\sigma_{\mathrm{tot}}^2=\sigma^2+\sigma_a^2$. At positive base noise, the derivative with respect to $\sigma_a$ is zero at $\sigma_a=0$ and positive at positive $\sigma_a$. That hypothetical mathematical statement is included in Section~A.3. It is deliberately not used to fill the missing implementation rule for the recorded \texttt{added\_sigma} field.

\subsection{Reading the Actual Probe Results}

The informative small-cut average is negative, but this is not because it has a large future penalty. The immediate/full ratio is $99.89827997\%$, and for the larger cut it is $99.98971257\%$. These are ratios of mean gains, not averages of individual percentage ratios. The remaining continuation differences are small relative to the loss at impact. Near-Nash baseline prices make a static margin-loss explanation consistent with the observed decomposition.

Consider informative seed 20261001, firm 0, time 64, cut 0.1. Its immediate mean gain is approximately $0.0006133863$ and its full mean gain is $0.0005733491$. The continuation term is about $-0.0000400371$: later changes reduce the benefit, but do not remove it. This probe is a concrete example of a profitable price cut despite a small negative continuation effect.

For informative seed 20261005, firm 0, time 64, cut 0.1, the full mean gain is $0.0038461616$, slightly above its immediate gain of about $0.003834$. The continuation effect does not constitute a loss in that example. Another positive informative probe occurs for the same seed and firm at time 256. The single positive replaced-channel small-cut probe occurs at seed 20261008, firm 0, time 256, with full gain about $0.0001222647$ and no meaningful continuation difference.

The corrected informative mean signed response area is $-1.3628785\times10^{-5}$ price-periods. Its seed-bootstrap interval excludes zero, so there is evidence of a small downward mean response. However, an area has different units from a discounted profit loss, and statistical detectability is not a criterion for an incentive constraint. The audit's conclusion follows from combining the small response with the negligible continuation losses and the existence of positive tested gains. It is not based on asserting that every informative response is zero.

The positive-probe intervals are unadjusted and conditional on particular trained policies. We report the four examples to expose heterogeneity, not as four independent population-level discoveries. The complete tables allow readers to inspect all cases rather than only an average or selected favorable trajectories.

\subsection{Welfare: A Correct Definition and a Separate Archived Field}\label{app:welfare}

For the stated logit demand, normalized expected consumer surplus is the inclusive value
\[
\mathrm{CS}(p_1,p_2)=\mu\log\!\left(1+e^{(v-p_1)/\mu}+e^{(v-p_2)/\mu}\right).
\]
Normalizing relative to the outside-option baseline removes the irrelevant additive utility constant. Standard total surplus is
\begin{equation}
\mathrm{TS}(p_1,p_2)=\mathrm{CS}(p_1,p_2)
+\pi_1(p_1,p_2)+\pi_2(p_2,p_1).
\label{eq:supp-welfare}
\end{equation}
It includes preference surplus, so the bound $\mathrm{TS}\le v-c$ does not follow. In particular, at the continuous symmetric Nash price the formula is about $1.1609510$, while $(v-c)(q_1+q_2)$ is about $0.9427536$. The simpler expression is not the same welfare concept.

The archive labels its output \texttt{mean\_total\_surplus}, with means near $1.1596$. These values are incompatible with an everywhere-bounded-by-one formula, but compatibility with a conventional formula is not proof of implementation identity. The simulator code and period-level price trajectories are missing. Because the welfare expression is nonlinear, substituting an aggregate mean price cannot reconstruct the mean welfare exactly either. We therefore retain the saved values under the explicit label ``archived welfare index'' and avoid making a verified economic welfare claim from them.

This distinction preserves both the mathematical definition and the empirical provenance boundary. Replacing a wrong definition by a plausible one must not silently certify unavailable code. Resolving the remaining question requires the original simulator implementation or sufficiently detailed period-level outputs.

\subsection{How Far the Evidence Can Be Carried}

The analytical model establishes a conditional incentive calculation for a fixed infinite-horizon public strategy. The observations establish properties of a finite set of deployed PPO policies under a fixed environment and finite audit family. The postprocessing establishes consistency of stored outputs, printed tables, bootstrap summaries, and recorded checkpoint hashes. Each layer has a different logical scope.

In particular, a known finite terminal date can change incentives and learning even though the discount is close to one. A 32-report actor cannot be treated as an arbitrary infinite-memory public strategy. A replacement that changes marginal report statistics cannot be interpreted as a pure information treatment. The zero-response negative control does not measure sensitivity to a known positive control. These limits remain after every reported calculation is reproduced correctly.

The current repository supports archive postprocessing. It does not contain the original \texttt{acrl} implementation, native simulator, posterior draws, calibration observations, or individual evaluation-rollout arrays. Accordingly, the revision neither claims a new training replication nor invents missing parameters. Section~G records the available identities, Section~H describes conditional reproduction requirements, and the final repository section lists the public artifacts and rechecking commands.


%% ============================================================
%% References (including works cited only in the supplement)
%% ============================================================

\bibliographystyle{iclr2027_conference}
\begin{thebibliography}{9}

\bibitem[Abreu et~al.(1990)]{abreu1990toward}
D.~Abreu, D.~Pearce, and E.~Stacchetti.
\newblock Toward a theory of discounted repeated games with imperfect
  monitoring.
\newblock \emph{Econometrica}, 58(5):1041--1063, 1990.

\bibitem[Calvano et~al.(2020)]{calvano2020artificial}
E.~Calvano, G.~Calzolari, V.~Denicol\`o, and S.~Pastorello.
\newblock Artificial intelligence, algorithmic pricing, and collusion.
\newblock \emph{American Economic Review}, 110(10):3267--3297, 2020.

\bibitem[Calvano et~al.(2021)]{calvano2021algorithmic}
E.~Calvano, G.~Calzolari, V.~Denicol\`o, and S.~Pastorello.
\newblock Algorithmic collusion with imperfect monitoring.
\newblock \emph{International Journal of Industrial Organization}, 79:102712,
  2021.

\bibitem[Cho et~al.(2014)]{cho2014gru}
K.~Cho, B.~van Merri\"enboer, C.~Gulcehre, D.~Bahdanau, F.~Bougares, H.~Schwenk, and Y.~Bengio.
\newblock Learning phrase representations using RNN encoder--decoder for statistical machine translation.
\newblock \emph{arXiv preprint arXiv:1406.1078}, 2014.

\bibitem[Deng et~al.(2025)]{deng2025exploring}
S.~Deng, M.~Schiffer, and M.~Bichler.
\newblock Exploring competitive and collusive behaviors in algorithmic
  pricing with deep reinforcement learning.
\newblock \emph{arXiv preprint arXiv:2503.11270}, 2025.

\bibitem[Fudenberg et~al.(1994)]{fudenberg1994folk}
D.~Fudenberg, D.~K.~Levine, and E.~Maskin.
\newblock The folk theorem with imperfect public information.
\newblock \emph{Econometrica}, 62(5):997--1039, 1994.

\bibitem[Green and Porter(1984)]{green1984noncooperative}
E.~J.~Green and R.~H.~Porter.
\newblock Noncooperative collusion under imperfect price information.
\newblock \emph{Econometrica}, 52(1):87--100, 1984.

\bibitem[Schulman et~al.(2017)]{schulman2017proximal}
J.~Schulman, F.~Wolski, P.~Dhariwal, A.~Radford, and O.~Klimov.
\newblock Proximal policy optimization algorithms.
\newblock \emph{arXiv preprint arXiv:1707.06347}, 2017.

\bibitem[Vehtari et~al.(2021)]{vehtari2021rank}
A.~Vehtari, A.~Gelman, D.~Simpson, B.~Carpenter, and P.-C.~B\"urkner.
\newblock Rank-normalization, folding, and localization: An improved
  $\hat{R}$ for assessing convergence of {MCMC}.
\newblock \emph{Bayesian Analysis}, 16(2):667--718, 2021.

\end{thebibliography}

\clearpage
\section{Repository and Available Artifacts}\label{app:repository}

The public repository contains the current main paper, this supplement, the recorded experiment archive, and the scripts used to recheck tables and regenerate figures. The source files are \texttt{main.tex} and \texttt{supplement.tex}; their compiled public versions are \texttt{main.pdf} and \texttt{supplement.pdf}. There is one current pair of manuscripts in the repository file tree. Earlier manuscript files and ZIP packages have been removed so that only the current public release remains.

To recheck the stored outputs, run \texttt{python3 scripts/verify\_archive.py} with \texttt{-\/-root .} and \texttt{-\/-output verification/archive-results.json}. The arithmetic checks use \texttt{python3 scripts/verify\_theory.py}; figures are regenerated by \texttt{python3 scripts/render\_figures.py}. These scripts postprocess archived summaries and check file identities. They do not replace the missing original simulator, training implementation, posterior/calibration inputs, or individual-rollout data, and do not deserialize checkpoints or train additional agents.

This public-release repository address identifies its owner; an anonymous review submission should use an anonymized artifact link instead. The public repository is:

\url{https://github.com/duomo77/beyond-high-prices-iclr2027}

\end{document}
